\section{Appendix: Liturgical Resolution}
\label{sec:liturgical-resolution}

{\small
\textbf{Celebration.} The Twenty-eighth Sunday in Ordinary Time, a Sunday in
Ordinary Time at place~6 of the table of liturgical days; green; Gloria and
Creed. The formula is \bid{PC-S54-A}, the Year~A target of the parent
\bid{PC-S54}; it has the single Sunday form. A Mass anticipated on Saturday
evening uses the same formulary, and there is no Vigil form.

\textbf{Books, calendar and territory.} The \work{Roman Missal}, Third
Edition, for use in the dioceses of the United States of America (2011), whose
Latin base is the \work{Missale Romanum}, editio typica tertia (2002), as
reprinted with emendations in 2008; the \work{Lectionary for Mass for Use in
the Dioceses of the United States of America}, second typical edition,
volume~I, no.~142, which follows the \work{Ordo lectionum Missae}, editio
typica altera (1981); the \work{General Instruction of the Roman Missal} in the
United States edition of 2011 with the emendations of 2021; the
\work{Universal Norms on the Liturgical Year and the Calendar}; and the
\work{Liturgical Calendar for the Dioceses of the United States of America} for
2026. Only the national calendar is applied. No diocesan, religious, parish,
titular, dedication or patronal calendar was supplied, and none is assumed.

\textbf{Cycle and occurrence.} Year~A of the Sunday cycle runs from 30~November
2025 to 22~November 2026, and the national calendar enters Sunday 11~October
2026 as the Twenty-eighth Sunday in Ordinary Time, green, with the readings of
Lectionary no.~142 and Psalter Week~IV. No other celebration stands on that
date in the calendar's dated entries. The weekday cycle governs no Sunday.

\textbf{The Missal formulary and the three cycles.} The antiphons and orations
belong to the Twenty-eighth Week in Ordinary Time and are common to Years A, B
and~C; their text is recorded once, in the collection's record for that week.
The readings, psalm and Alleluia verse are those of Year~A alone. The second
reading is the last of four in the Lectionary's course through Philippians in
Year~A, and the Gospel belongs to its semi-continuous reading of Matthew,
between the parable of the vineyard on the previous Sunday and the question of
tribute on the next.

\textbf{Branches.} No selection among the following was supplied, and none is
inferred. Both forms of the Gospel and both Communion antiphons are studied in
full; a claim resting on Matthew 22:11--14 does not hold for a celebration
that uses the shorter form, and a claim resting on one Communion antiphon holds
only where it is chosen.
\par}

\begin{branchtable}
\bid{gospel-long-form} & Lectionary no.~142 prints Mt 22:1--14 as the longer
form & appointed alternative & Gospel & Unselected; studied in full. The
choice follows the hearers' capacity to listen with profit (\work{Praenotanda}
80).\\
\bid{gospel-short-form} & Lectionary no.~142 prints Mt 22:1--10 as the shorter
form & appointed alternative & Gospel & Unselected; studied in full. It omits
the guest without a wedding garment and the saying of v.~14.\\
\bid{communion-antiphon-psalm} & The Missal prints cf.\ Ps 34 (33):11 first &
appointed alternative & Communion Antiphon & Unselected; studied in full.\\
\bid{communion-antiphon-epistle} & The Missal prints 1~Jn 3:2 second, under
\latin{Vel} & appointed alternative & Communion Antiphon & Unselected; studied
in full. Rubric~6 for Ordinary Time prefers an antiphon that agrees with the
Gospel; neither is drawn from a Gospel, and the books do not say whether either
agrees with this one.\\
\bid{entrance-or-communion-chant-substitution} & \work{General Instruction} 48
and 87: the Missal or \work{Graduale Romanum} antiphon, the \work{Graduale
Simplex}, a chant from another approved collection of psalms and antiphons, or
another suitable liturgical chant, the last two approved by the Conference of
Bishops or the Diocesan Bishop & permitted & Entrance; Communion & Unselected;
the Missal antiphons are the texts studied.\\
\bid{psalm-or-acclamation-chant-substitution} & \work{General Instruction} 61
and 62: a seasonal common psalm when the psalm is sung; the \work{Graduale}
forms or another approved collection of psalms and antiphons, but never a song
or hymn in place of the psalm & permitted & Responsorial Psalm; Alleluia verse
& Unselected; the Lectionary's psalm and verse are the texts studied.\\
\bid{offertory-chant} & \work{General Instruction} 74 & permitted & Preparation
of the gifts & Unselected; the Missal prints no Offertory antiphon.\\
\bid{sprinkling-rite} & \work{General Instruction} 51 and the Order of Mass:
on Sundays the blessing and sprinkling of water may replace the Penitential
Act & ritual substitution & Introductory Rites & Unselected; it brings its own
texts.\\
\bid{compatible-preface-and-eucharistic-prayer} & Rubric~5 for Ordinary Time;
\work{General Instruction} 364--365 & permitted & Eucharistic Prayer &
Unselected; one of the eight Prefaces of the Sundays in Ordinary Time with a
compatible Eucharistic Prayer, as one coupled choice.\\
\bid{eucharistic-prayer-iv-with-preface} & \work{General Instruction} 365~d &
conditional & Eucharistic Prayer & Unselected; Eucharistic Prayer~IV keeps its
own invariable Preface.\\
\bid{concluding-blessing-form} & The Missal's solemn blessings and prayers over
the people, at the priest's discretion & permitted & Concluding Rites &
Unselected; none is appointed to this Sunday.\\
\bid{local-proper-calendar} & \work{Universal Norms} 59--60 with a particular
calendar & conditional & Whole Mass & Not supplied; a local celebration would
need its own resolution.\\
\end{branchtable}
